\section{SocCogEval: Task, Protocol, and Measurement}\label{sec:method}\label{sec:task}

SocCogEval combines a task definition, a factorial protocol, and a measurement suite. The task defines what would count as social cognitive transfer. The protocol separates interaction from briefs, initialization noise, and memory machinery. The measurements convert each run into contrast-level evidence.

\subsection{Task: Decision, Prediction, Transfer}\label{sec:sct-task}

An SCT instance specifies a frozen base model $M$, an agent population $A{=}\{a_1, \dots, a_k\}$ of $k$ clone agents that share system prompt and decoding parameters, an interaction graph $G$ over $A$ (all-to-all, two-component segregated, or empty), a scenario $S$ with a 5-dim allocation space, an optional brief assignment $T$ (asymmetric A and B, symmetric A, or none), and an interaction horizon $N$ (default 30 rounds). The clone constraint is the load-bearing one. In any condition without briefs, $a_i$ and $a_j$ have bit-identical initial state. When briefs are assigned, the only initial difference is the private brief document. Three nested questions follow. SCT-Decision asks whether post-interaction allocation in segregated, asymmetrically briefed conditions diverges beyond what brief alone or initialization noise produce. SCT-Prediction asks whether agents in segregated conditions predict in-group members more accurately than out-group members. SCT-Transfer asks whether the divergence appears on instruments topically unrelated to the discussion, after subtracting the no-interaction baseline.

\subsection{Factorial Social Interaction Protocol}\label{sec:fsip}

The factorial axes are brief ($T \in \{\text{none}, \text{A only}, \text{asymmetric A and B}\}$) and interaction ($N \in \{0, 30\}$ with topology $G \in \{\text{open}, \text{segregated}\}$). The design uses five conditions plus one mechanism control (Figure~\ref{F2}). $C_1$ (open, symmetric brief) uses an all-to-all topology in which every agent receives Brief A, as an open-network reference. $C_2$ (segregated, asymmetric brief) splits agents into two components, half with Brief A and half with Brief B, and is the headline condition. $C_3$ (segregated, no brief) keeps the $C_2$ topology without briefs and isolates interaction alone. $C_4$ (brief only) assigns the A and B briefs with $N{=}0$ and isolates brief alone. $C_5$ (no brief, no interaction) sets the initialization noise floor. $C_6$ (solo generation) keeps the $C_2$ briefs, memory policy, probe cadence, and logical call budget, with private generation in place of partners (Figure~\ref{F3sch}).

The load-bearing causal contrast is $(C_2){-}(C_4)$ on every measured outcome. It holds the brief fixed, since both conditions use asymmetric A and B briefs, and toggles 30 rounds of interaction on or off. Interaction-created history, memory, and elapsed conversation are post-treatment pathways through which the treatment acts, so $(C_2){-}(C_4)$ is the total effect of interaction given the same brief in this protocol. One scope limit matters. $C_1$ uses symmetric briefs while $C_2$ uses asymmetric briefs, so the $(C_2){-}(C_1)$ contrast varies brief-symmetry and topology jointly and does not isolate topology. We therefore exclude topology from the causal claims and report $(C_2){-}(C_1)$ only as a descriptive comparison. Replica counts per condition are $\{C_1{=}7, C_2{=}8, C_3{=}5, C_4{=}5, C_5{=}3, C_6{=}5\}$ per scenario, with $k{=}8$ agents per replica.

\textbf{Solo-generation control.} $C_6$ gives each agent three private generation calls per round (initial reasoning, reconsideration, allocation reflection) and one memory extraction, with no partner identity, messages, or feedback. This yields 720 content and 240 extraction calls per run, matching $C_2$'s intended logical totals. Its design was frozen before formal outcomes were examined, and all ten formal $C_6$ runs passed audit (Appendix~\ref{app:c6}). $(C_6){-}(C_4)$ measures the solo-generation shift, and $(C_2){-}(C_6)$ compares social and solo histories.

\subsection{Scenarios, Runtime Cycle, and Models}\label{sec:fsip-cycle}

We run two topically disjoint scenarios. Greenfield is a community-grant allocation problem with \$10M across five categories such as Education, Healthcare, and Infrastructure. Riverton is a workforce transition budget problem with \$5M across categories such as Automation and Workforce Training. Both use 5-dim allocation vectors so Part 1 metrics are commensurable. Both have softened A and B briefs that frame the same problem through moderately divergent value lenses (Appendix~\ref{app:scenarios}). Per-round mechanics follow a Build-Talk-Probe-Persist cycle. Build updates a typed Social-Relational Memory (SRM) hypergraph from the previous round's conversations. Talk forms 4 random pair conversations among 8 agents with 6 turns each. Probe runs SocCogEval instruments at predetermined cadences, with belief and common-ground probes every five rounds that leave conversation history and memory unchanged. Persist writes the round artifacts to disk.

\textbf{Models and filtering.} The primary model is Qwen3-122B-FP8~\citep{Qwen2025_Qwen3} served locally via vLLM. Five model variants run on $\{C_2, C_4, C_5\}$ in both scenarios, three runs per condition and scenario: Qwen3-VL-32B-Instruct~\citep{Qwen2025_Qwen3VL} run text-only, DeepSeek-V3~\citep{DeepSeekAI2024_DeepSeekV3}\footnote{Historical DeepSeek $C_2$ runs used the provider identifier \texttt{deepseek-chat} and the added $C_4$ runs used \texttt{deepseek-v4-flash}. Exact provider revisions were not recorded (Appendix~\ref{app:routing}).}, GPT-4.1-mini~\citep{OpenAI2025_GPT41}, Gemini-2.5-Flash~\citep{GoogleDeepMind2025_Gemini25}, and Llama-4-Scout~\citep{MetaAI2025_Llama4}. Figures abbreviate the six models as Q-122B, Q-VL, DS-V3, GPT-mini, Gem-Flash, and L4-Scout. The selected dataset contains 168 runs: 56 primary, 90 cross-variant, 12 no-SRM, and 10 $C_6$. Primary factorial analyses use a strict run-id filter so cross-variant replicas cannot enter the within-Qwen3-122B-FP8 baseline. Content generation uses temperature 0.7, BFI-10 and conformity use 0.7, and allocation, prediction, and memory extraction use 0.

\subsection{Measurements and Statistics}\label{sec:metrics}

Among the Part 2 instruments, we designate BFI-10 response-profile drift as the primary transfer outcome. It is topic-independent, compact, and inherited from classical BFI-10 psychometrics and prior LLM psychometric evaluation work~\citep{Rammstedt2007_BFI10,PsychoBench,SerapioGarcia2025_PsychometricLLM,Ye2025_LLM_Psychometrics}. BFI-10 is manuscript-designated primary, without an external timestamped preregistration. The claim does not require the inventory to be psychometrically valid for LLMs in the human sense. We use its fixed item structure as an off-topic response-profile probe whose items are unrelated to either allocation scenario. Conformity and cross-prompt consistency are secondary transfer instruments, using Asch-style conformity materials and CogMir-inspired probe phrasing~\citep{Asch1956_Conformity,CogMir}. Allocation divergence and prediction MAE are the Part 1 instruments and support the Decision and Prediction parts of the SCT route (scoring formulas in Appendix~\ref{app:stats}).

For agent $i$ in run $r$, let $s_{rif}^{t}$ be the score of BFI factor $f$ at phase $t\in\{\mathrm{pre},\mathrm{post}\}$, after reverse-scoring items 1, 3, 4, 5, and 7. Agent drift, run-mean drift, the contrast between conditions, and the signed factor contrast are
\begin{equation}\label{eq:drift}
 d_{ri}=\sum_{f=1}^{5}\bigl|s_{rif}^{\mathrm{post}}-s_{rif}^{\mathrm{pre}}\bigr|,\quad
 D_r=\tfrac{1}{8}\sum_{i}d_{ri},\quad
 \bar\Delta_{a,b}=\overline D_{C_a}-\overline D_{C_b},\quad
 \kappa_f=\overline m_{C_2,f}-\overline m_{C_4,f},
\end{equation}
where $m_{rf}=\frac{1}{8}\sum_i(s_{rif}^{\mathrm{post}}-s_{rif}^{\mathrm{pre}})$. Because $d_{ri}$ sums five factor-level shifts, $\bar\Delta{=}0.92$ corresponds to about 0.18 Likert points per factor. Unsigned L1 can grow through directional movement or dispersion, so $\kappa_f$ and run-level variance ratios report direction and spread separately.

\textbf{Unit of analysis.} Agents in one run share an evolving history and interfere with each other's outcomes~\citep{Hudgens2008,Aronow2017}, so the run is the independent unit for the primary inference and agents within a run are not replicates~\citep{Hurlbert1984}. Each run is summarized by $D_r$. Conditions are compared with two-sided exact permutation tests that enumerate all label allocations, so no Monte Carlo correction is needed~\citep{Phipson2010}, and with run-bootstrap 95\% intervals~\citep{Efron1979}. Raw $\bar\Delta$ leads the reporting. Hedges' $g$ with the small-sample correction~\citep{Hedges1981} is secondary, and we report it as undefined when the pooled within-condition SD falls below 0.05 in the metric's native units. The gate status in Table~\ref{tab:gates} keeps the original agent-stratified two-level cluster bootstrap, which resamples replicas and then agents within each replica (10{,}000 resamples), with Welch's $t$ as the parametric reference.

\textbf{Reporting standard.} The primary SCT-positive claim is made for BFI-10 response-profile drift. It requires three criteria. First, a Part 1 allocation contrast must clear the majority-of-dimensions rule on $(C_2){-}(C_5)$. Second, BFI-10 must clear its CI gate on $(C_2){-}(C_5)$. Third, BFI-10 must clear its CI gate on the pure-interaction $(C_2){-}(C_4)$. Secondary instruments get the same gate status but serve as scenario-specific supporting evidence, not the load-bearing claim.
